\subsection{出次数1の行列式非零部分の全排除}
\label{subsec:cubic-CY-four-point-two-open}
\begin{theorem}\label{thm:cubic-CY-four-point-two-open-empty}
四単純基点の元表示で
\[
 x_0x_1\ne0,\qquad x_3=x_4=0,\qquad x_0x_2-x_1^2\ne0
\]
を満たす全領域と、その全monomial軌道はH3包絡を与えない。
$x_5=0$ と $x_5\ne0$ の両方を含む。
\end{theorem}
\begin{proof}
固定出典\cite{AS4CYpinned}の第二共有開部分と第五係数零開部分の必要被覆を用いる。
両方の全controllerを隔離出力先で再実行し、当repoでも以下の全元係数消去・
全局所化分岐・表示移送を独立に確認する。
AS正則性の十分性や全同型群は用いない。
まず独立な対角比で $x_0=x_1=1$ とし、$p=x_2\ne1$ と書く。

\paragraph{$x_5\ne0$}
さらに $x_5=1$ とする。
全21係数のSchur式から $\lambda,\lambda+1,p-1$ だけを単元として消去する。
$q=x_{10},u=x_6,v=x_7$ を残すと、元残余idealは
\[
 q(u-v),\qquad q(pu-v)
\]
を含む。$q\ne0$ なら $p\ne1$ により $u=v=0$。
これを元残余式へ代入したidealは $q\lambda(\lambda+1)(p-1)$ の局所化で単位idealになる。
従って $q=0$ であり、この側の全単元消去は残余なしで完了する。
全係数は既存 $H_{\lambda,p}$、すなわち
\[
\begin{gathered}
 x_0=x_1=x_5=1,\quad x_2=p,\quad x_9=(\lambda^2-1)/16,\\
 x_{14}=\lambda^2/(16(p-1)),\quad x_{19}=-\lambda/(16(p-1)),\quad
 x_{20}=p/(16(p-1)),\quad g=1/8
\end{gathered}
\]
となり、他は零である。
独立再現器は全消去式・単元pivot・最終全係数を検査し、
元関係が既存表示と文字どおり一致することを確認した。
命題\ref{prop:cubic-second-shared-H}から全許容値でH3でない。

\paragraph{$x_5=0$}
同じ局所化環で全21係数を消去する。
$q=0$ の全残余は $\lambda=1$ を単元倍の式として強制し、
他の全係数は
\[
\begin{gathered}
 x_0=x_1=1,\quad x_2=p,\quad x_9=r,\quad
 x_{14}=1/(16(p-1)),\\
 x_{19}=-1/(16(p-1)),\quad x_{20}=p/(16(p-1)),\quad g=1/8
\end{gathered}
\]
となる。$r$ を自由に保って元全係数へ戻ることも検査した。
対角 $\operatorname{diag}(1,z,z,1)$ は $r$ を $z^2r$ へ移すので、
$r=0$ 又は1へ全規格化でき、既存 $M^0_p,M^1_p$ に尽くされる。
命題\ref{prop:cubic-zero-fifth-m}で全排除される。

残る $q\ne0$ で $p=0$ の場合には、
\[
 P=\operatorname{diag}(1,1/\lambda,1/\lambda,1/\lambda)P_{(bc)}
\]
を施すと、元係数のまま $x'_0=x'_3=x'_4=0,x'_1=1$ となる。
定理\ref{thm:cubic-CY-four-point-first-empty}がこの全chartを排除するので、
正則性の十分条件や特定の残余係数の零を追加してはいない。

以後 $pq\ne0$。
元残余式の一つが単元倍を除き
\[
 p(-\lambda^2+16qx_8)(-\lambda q+\lambda x_8-q)=0
\]
となる。二因子をそれぞれ零とした全必要分岐は
\[
\begin{array}{c|c|c|c}
 &x_8&(\lambda,p)&q^2\\\hline
 A&\lambda^2/(16q)&(1/2,-2)&-1/192\\
 B&(\lambda+1)q/\lambda&(2,-2)&-1/12
\end{array}
\]
である。
これは一般点の数値探索ではなく、$x_6,x_7,x_9$ を含まない元残余式のidealと、
表の三式のidealとが、$\lambda(\lambda+1)p(p-1)q$ の局所化で一致することによる。
両方のidealを有理数体上で独立に再計算した。
各行の $q$ の一根を代入した残余は、数体上のidealとして
$x_6,x_7,x_9$ の各正冪を含むので、全て零となる。
他の根は数体共役であり、同じ結論を持つ。
ここでの幾何学的な必要被覆を、残余schemeの被約性と同一視していない。

この四候補はいずれも既存の固定表示 $N$ と次数付き同型である。
当repoでは全24置換と対角正規化を元関係へ作用させ、
四候補の各々に対して可逆な関係基底移送を具体的に確認した。
$N$ の元表示も既存排除表示と文字どおり一致するので、
命題\ref{prop:cubic-zero-fifth-fixed}から全てH3でない。
以上で二開chartの全場合が尽くされた。
\end{proof}

独立検証器は\\
\texttt{verify\_cubic\_CY\_four\_point\_two\_open.py}。
最小固定入力は必要な消去手順・残余式・候補係数だけを保持する。
行列式零境界、出次数2・部分活性・全共有支持の全必要係数被覆と相対分類は未完成である。
